% Generated draft: requires mathematical and source review.
\documentclass[11pt]{article}
\usepackage[margin=25mm]{geometry}
\usepackage{fontspec}
\setmainfont{Latin Modern Roman}
\usepackage{amsmath,amssymb,mathtools,unicode-math}
\setmathfont{Latin Modern Math}
\usepackage{enumitem,needspace,xurl,hyperref}
\setlength{\parindent}{0pt}
\setlength{\parskip}{5pt}
\setlength{\emergencystretch}{3em}
\newcommand{\N}{\mathbb{N}}
\newcommand{\Z}{\mathbb{Z}}
\newcommand{\Q}{\mathbb{Q}}
\newcommand{\R}{\mathbb{R}}
\newcommand{\C}{\mathbb{C}}
\newcommand{\F}{\mathbb{F}}
\newcommand{\A}{\mathbb{A}}
\newcommand{\PP}{\mathbb{P}}
\newcommand{\E}{\mathbb{E}}
\newcommand{\eps}{\varepsilon}
\newcommand{\dd}{\,\mathrm{d}}
\begin{document}
\section*{UnsolvedMath ID: 700540}
\textbf{Source problem code:} catalogue entry
\subsection{Definitions and mathematical statement}
Work in the language of groups, with multiplication, inversion, and the identity element. A group identity is a universally quantified equation between group terms. A variety of groups is the class of all groups satisfying some set of group identities.

A group quasi-identity is a universally quantified implication of the form
\[
\forall x_1,\ldots,x_n\quad
\left(\bigwedge_{i=1}^{m} u_i(x_1,\ldots,x_n)=v_i(x_1,\ldots,x_n)\right)
\Longrightarrow u(x_1,\ldots,x_n)=v(x_1,\ldots,x_n),
\]
where all displayed expressions are group terms and the conjunction is finite; the case \(m=0\) is allowed. A quasivariety of groups is the class of all groups satisfying some set of group quasi-identities. A subquasivariety of a variety \(\mathcal V\) is a quasivariety \(\mathcal Q\) with \(\mathcal Q\subseteq\mathcal V\).

A variety \(\mathcal V\) is called \emph{primitive} if every subquasivariety \(\mathcal Q\subseteq\mathcal V\) is itself a variety of groups. A group is \emph{locally finite} if every subgroup generated by finitely many elements is finite. A variety is locally finite if each of its groups is locally finite.

\textbf{Problem (V. A. Gorbunov, 1984; Kourovka Notebook 9.5).} Describe all primitive varieties of groups. In particular, determine whether every primitive variety of groups is locally finite; equivalently, whether
\[
\forall\mathcal V\;\left[\mathcal V\text{ is a primitive variety of groups}\ \Longrightarrow\ 
\forall G\in\mathcal V\;\forall S\subseteq G\;
\bigl(|S|<\infty\Longrightarrow |\langle S\rangle|<\infty\bigr)\right].
\]
The supplied registry labels this problem open; that label is supplied source metadata and is not an independently verified assessment.
\subsection{Short English statement}
Classify the varieties of groups for which every contained quasivariety is a variety. Determine also whether every such variety consists entirely of groups whose finitely generated subgroups are finite.
\subsection{Sources}
\begin{itemize}
\item Supplied local UnsolvedMath registry record, permanent ID \(700540\), exact source code \texttt{catalogue entry}, title \emph{Kourovka Notebook Problem 9.5}; upstream revision \texttt{d371ee917ac9647fbb06cd6fc529deec6f3dda82}. The record supplies the status \emph{open}.
\item \emph{The Kourovka Notebook (Unsolved Problems in Group Theory)}, Problem 9.5, proposed by V. A. Gorbunov in the 9th issue (1984). The supplied background also identifies the 21st issue edition (2026), \texttt{arXiv:1401.0300v48}. The supplied collection description names E. I. Khukhro and V. D. Mazurov as editors. These bibliographic details are supplied attribution and have not been independently checked.
\end{itemize}
\end{document}
